\documentclass[10pt]{beamer}

% ===================== THEME =====================
\usetheme{Madrid}
\usecolortheme{whale}
\setbeamertemplate{navigation symbols}{}
\setbeamertemplate{footline}[frame number]
\setbeamercovered{transparent}

% ===================== LANGUE / ENCODAGE =====================
\usepackage[french]{babel}
\usepackage[utf8]{inputenc}
\usepackage[T1]{fontenc}

% ===================== PACKAGES =====================
\usepackage{booktabs}
\usepackage{array}
\usepackage{tikz}
\usetikzlibrary{shapes.geometric,arrows.meta,positioning,fit,backgrounds,decorations.pathreplacing}
\usepackage{amsmath}
\usepackage{amssymb}
\usepackage{multicol}
\usepackage{tabularx}

\definecolor{institbleu}{RGB}{0,32,91}
\definecolor{institor}{RGB}{193,154,58}
\setbeamercolor{structure}{fg=institbleu}
\setbeamercolor{title}{fg=white,bg=institbleu}
\setbeamercolor{frametitle}{fg=white,bg=institbleu}

% ===================== PAGE DE TITRE =====================
\title[Cartographie \& mise en relation des entreprises]{Cartographie des besoins, conception et réalisation \\ d'une plateforme de mise en relation des entreprises}
\subtitle{Soutenance en vue de l'obtention du diplôme d'Ingénieur en Data Science}
\author[Landry]{Présenté par : \textbf{Landry}\\ \small Matricule : 21D0180EP}
\institute[ENSPM]{
	\'Ecole Nationale Supérieure Polytechnique de Maroua \\
	Université de Maroua \\[4pt]
	Département Informatique et Télécommunication / Data Science
}
\date{\today}

\begin{document}
	
	% ---------------------------------------------------------------
	% PARTIE 1 : PAGE DE TITRE ET CONTEXTE INSTITUTIONNEL
	% ---------------------------------------------------------------
	\begin{frame}
		\titlepage
	\end{frame}
	
	\begin{frame}{Cadre de la soutenance}
		\begin{columns}[T]
			\begin{column}{0.55\textwidth}
				\textbf{Encadrement académique} \\
				Pr. Kaladzafi Guidedi \\[8pt]
				\textbf{Encadrement professionnel} \\
				M. Abdoulaziz Jumnbe \\[8pt]
				\textbf{Structure d'accueil}\\
				Ministère de l'Économie, de la Planification et de l'Aménagement du Territoire (MINEPAT)\\
				Direction Générale de l'Économie et de la Planification des Investissements Publics (DGEPIP)\\
				Division des Analyses et des Politiques Économiques (DAPE)
			\end{column}
			\begin{column}{0.45\textwidth}
				\begin{tikzpicture}[scale=0.9]
					\node[draw, rounded corners, fill=institbleu!10, minimum width=3.2cm, minimum height=1cm, align=center] (m) at (0,2) {\small MINEPAT};
					\node[draw, rounded corners, fill=institbleu!15, minimum width=3.4cm, minimum height=1cm, align=center] (dg) at (0,0.7) {\small DGEPIP};
					\node[draw, rounded corners, fill=institbleu!25, minimum width=3.4cm, minimum height=1cm, align=center] (dape) at (0,-0.6) {\small DAPE};
					\draw[-{Latex}] (m) -- (dg);
					\draw[-{Latex}] (dg) -- (dape);
				\end{tikzpicture}
			\end{column}
		\end{columns}
	\end{frame}
	
	\begin{frame}{Plan de la présentation}
		\small
		\begin{multicols}{2}
			\tableofcontents
		\end{multicols}
	\end{frame}
	
	% ---------------------------------------------------------------
	% PARTIE 2 : CONTEXTE
	% ---------------------------------------------------------------
	\section{Contexte}
	
	\begin{frame}{Contexte de l'étude (1/2)}
		\begin{itemize}
			\item Dans de nombreux pays, le \textbf{secteur privé} joue un rôle déterminant dans la production des biens et services ainsi que dans la création de valeur ajoutée.
			\item À travers les \textbf{trois piliers de la SND30}, le Cameroun se fixe comme objectif le développement du secteur privé comme moteur de la croissance économique.
		\end{itemize}
		\begin{center}
			\begin{tikzpicture}[scale=0.85, every node/.style={align=center, font=\scriptsize}]
				\node[draw, fill=institbleu, text=white, rounded corners, minimum width=3.1cm, minimum height=1.4cm] (p1) at (-3.4,0) {Transformation\\structurelle};
				\node[draw, fill=institor, text=white, rounded corners, minimum width=3.1cm, minimum height=1.4cm] (p2) at (0,0) {Développement\\du capital humain};
				\node[draw, fill=institbleu!70, text=white, rounded corners, minimum width=3.1cm, minimum height=1.4cm] (p3) at (3.4,0) {Gouvernance,\\décentralisation\\\& modernisation};
				\node[font=\bfseries] at (0,1.1) {SND30 -- 3 piliers};
			\end{tikzpicture}
		\end{center}
	\end{frame}
	
	\begin{frame}{Contexte de l'étude (2/2)}
		\begin{itemize}
			\item Des \textbf{réformes} ont été engagées et des \textbf{indicateurs de croissance} définis pour suivre la mise en \oe uvre de la SND30.
			\item L'\textbf{évaluation à mi-parcours} de l'atteinte de ces indicateurs fait état d'une économie camerounaise \textbf{résiliente}, mais encore \textbf{loin des idéaux fixés}.
		\end{itemize}
		\vspace{0.3cm}
		\begin{center}
			\begin{tikzpicture}[scale=0.95]
				\draw[->] (0,0) -- (6.5,0) node[right, font=\scriptsize] {Horizon SND30};
				\draw[->] (0,0) -- (0,3) node[above, font=\scriptsize] {Indicateurs};
				\draw[dashed, institbleu] (0,2.6) -- (6.2,2.6) node[right, font=\scriptsize, black] {Cible};
				\draw[thick, institor] (0,0.3) .. controls (2,0.9) and (3.5,1.3) .. (6,1.6);
				\node[font=\scriptsize] at (3,1.9) {Trajectoire observée (mi-parcours)};
			\end{tikzpicture}
		\end{center}
	\end{frame}
	
	% ---------------------------------------------------------------
	% PARTIE 3 : PROBLEMATIQUE
	% ---------------------------------------------------------------
	\section{Problématique}
	
	\begin{frame}{Problématique}
		\begin{itemize}
			\item Le secteur privé est incontournable pour la croissance économique d'un pays ; au Cameroun, ce secteur rencontre d'énormes difficultés, comme le montre l'évaluation à mi-parcours de la SND30.
			\item Permettre au secteur privé de contribuer davantage à la croissance suppose que l'administration centrale aide les entreprises à se défaire des besoins qui freinent leur développement.
			\item Plusieurs initiatives ont déjà été engagées dans cette optique, avec des résultats peu concluants \textit{(liste exhaustive des enquêtes et initiatives présentée dans le mémoire, chapitre 2)}.
			\item En 2025, le MINEPAT lance une \textbf{enquête de cartographie des besoins des entreprises}, avec pour objectifs principaux :
			\begin{itemize}
				\item permettre des \textbf{actions ciblées} de l'administration ;
				\item déboucher sur une \textbf{plateforme de collaboration} pour les entreprises.
			\end{itemize}
		\end{itemize}
	\end{frame}
	
	% ---------------------------------------------------------------
	% PARTIE 4 : DEFINITIONS ET CONCEPTS
	% ---------------------------------------------------------------
	\section{Définitions et concepts clés}
	
	\begin{frame}{Le besoin}
		\begin{columns}[T]
			\begin{column}{0.55\textwidth}
				\small
				D'après le \textit{Larousse}, un \textbf{besoin} désigne une exigence née d'un sentiment de manque ou de privation de quelque chose de nécessaire à la vie ou au bon fonctionnement d'un système.\\[6pt]
				Transposé à l'entreprise : ensemble des manques, insuffisances ou absences de ressources qui, une fois comblés, contribuent à la \textbf{croissance}, à la \textbf{compétitivité} et à la \textbf{pérennité} de l'entreprise.
			\end{column}
			\begin{column}{0.45\textwidth}
				\begin{tikzpicture}[scale=0.8]
					\draw[thick] (0,0) circle (1.3);
					\fill[institbleu!30] (0,0) circle (1.3);
					\draw[thick, dashed] (0,0) circle (0.75);
					\node[font=\scriptsize, align=center] at (0,0) {État\\actuel};
					\node[font=\scriptsize, align=center] at (0,-1.7) {Niveau requis};
					\draw[<->, thick, institor] (0.75,0) -- (1.3,0);
					\node[font=\scriptsize, above] at (1.0,0.15) {besoin};
				\end{tikzpicture}
			\end{column}
		\end{columns}
	\end{frame}
	
	\begin{frame}{La cartographie}
		\begin{columns}[T]
			\begin{column}{0.55\textwidth}
				\small
				Toujours selon le \textit{Larousse}, la \textbf{cartographie} est l'ensemble des opérations ayant pour objet l'élaboration, la rédaction et l'édition de cartes.\\[6pt]
				La \textbf{cartographie des besoins des entreprises} : ensemble des méthodes concourant à la représentation de la distribution des besoins des entreprises selon des variables identifiées.
			\end{column}
			\begin{column}{0.45\textwidth}
				\begin{tikzpicture}[scale=0.75]
					\draw[step=0.7, gray!40, thin] (0,0) grid (2.8,2.1);
					\fill[institor!70] (0.7,0.7) rectangle (1.4,1.4);
					\fill[institbleu!60] (1.4,1.4) rectangle (2.1,2.1);
					\fill[institbleu!30] (0,0) rectangle (0.7,0.7);
					\node[font=\tiny] at (1.4,-0.3) {Distribution spatiale des besoins};
				\end{tikzpicture}
			\end{column}
		\end{columns}
	\end{frame}
	
	\begin{frame}{La mise en relation}
		\small
		La \textbf{mise en relation des entreprises} renvoie au mécanisme par lequel deux ou plusieurs unités économiques sont amenées à se connaître, à échanger sur leurs capacités et leurs besoins respectifs, et à explorer les conditions d'une collaboration mutuellement bénéfique.
		\vspace{4pt}
		
		\alert{À distinguer du partenariat}, notion souvent confondue avec la mise en relation, mais qui implique un engagement contractuel formalisé entre les parties.
		\begin{center}
			\begin{tikzpicture}[scale=0.85]
				\node[draw, circle, fill=institbleu!30, minimum size=1.1cm] (a) at (0,0) {A};
				\node[draw, circle, fill=institor!40, minimum size=1.1cm] (b) at (3,0) {B};
				\draw[<->, thick] (a) -- (b) node[midway, above, font=\scriptsize] {échange besoins / offres};
			\end{tikzpicture}
		\end{center}
	\end{frame}
	
	% ---------------------------------------------------------------
	% PARTIE 5 : ETAT DE L'ART
	% ---------------------------------------------------------------
	\section{État de l'art}
	
	\begin{frame}{État de l'art -- Études de cartographie des besoins}
		\scriptsize
		\begin{tabularx}{\textwidth}{>{\raggedright\arraybackslash}p{2.6cm} X X}
			\toprule
			\textbf{Étude / initiative} & \textbf{Objectifs} & \textbf{Limites} \\
			\midrule
			Recensement Général des Entreprises (INS) & Dénombrer et caractériser le tissu économique national & Périodicité longue, faible granularité des besoins exprimés \\
			Enquête de climat des affaires (GICAM) & Mesurer la perception des entreprises sur l'environnement des affaires & Approche perceptive, peu actionnable individuellement \\
			Enquêtes sectorielles ministérielles antérieures & Identifier des contraintes par filière & Couverture partielle, absence de suivi dans le temps \\
			Enquête MINEPAT 2025 (cadre du mémoire) & Cartographier les besoins et alimenter une plateforme de mise en relation & En cours de valorisation -- objet de ce travail \\
			\bottomrule
		\end{tabularx}
	\end{frame}
	
	\begin{frame}{État de l'art -- Outils de mise en relation d'entreprises}
		\scriptsize
		\begin{tabularx}{\textwidth}{>{\raggedright\arraybackslash}p{2.2cm} X X}
			\toprule
			\textbf{Plateforme} & \textbf{Objectif principal} & \textbf{Limites} \\
			\midrule
			LinkedIn & Réseautage professionnel généraliste & Peu spécialisé pour l'appariement fin des besoins d'entreprises \\
			Kompass & Annuaire B2B mondial d'entreprises & Fonction essentiellement répertoire, appariement limité \\
			Alibaba & Mise en relation acheteurs-fournisseurs à l'international & Orienté commerce de marchandises, peu adapté au contexte local \\
			Go Africa Online & Annuaire d'entreprises africaines & Absence de moteur de matching basé sur les besoins exprimés \\
			\bottomrule
		\end{tabularx}
		\vspace{6pt}
		\normalsize
		\textbf{Constat} : aucun de ces outils ne combine \textit{cartographie statistique des besoins} et \textit{moteur de mise en relation ciblé}, adapté au contexte camerounais.
	\end{frame}
	
	% ---------------------------------------------------------------
	% PARTIE 6 : JUSTIFICATION DE L'ENQUETE ET DE LA PLATEFORME
	% ---------------------------------------------------------------
	\section{Justification de l'enquête et de la plateforme}
	
	\begin{frame}{Pourquoi cette enquête et cette plateforme ?}
		\begin{itemize}
			\item \textbf{Objectifs de l'enquête MINEPAT} : identifier finement les besoins des entreprises et permettre des actions ciblées.
			\item \textbf{Réponse aux limites} des enquêtes précédentes : granularité individuelle, actualisation continue, exploitation numérique directe des données.
			\item \textbf{Cohérence avec la SND30} : appui direct au pilier de transformation structurelle en levant les freins identifiés du secteur privé.
			\item \textbf{Opportunité de l'application} : transformer une base de données d'enquête en un \textbf{outil opérationnel} d'aide à la décision et de mise en relation.
			\item \textbf{Limites résolues} : absence d'outil dédié combinant statistique et appariement automatisé.
		\end{itemize}
	\end{frame}
	
	% ---------------------------------------------------------------
	% PARTIE 7 : PRESENTATION DES DONNEES
	% ---------------------------------------------------------------
	\section{Présentation des données}
	
	\begin{frame}{Les données de l'enquête}
		\small
		Données stockées sur un serveur sous deux tables, exportées au format \texttt{.csv} :
		\begin{itemize}
			\item \texttt{identifications.csv} : identification précise des entreprises.
			\item \texttt{besoins.csv} : besoins exprimés par les entreprises.
		\end{itemize}
		\vspace{6pt}
		\begin{columns}[T]
			\begin{column}{0.5\textwidth}
				\textbf{Éléments descriptifs (identifications)}
				\begin{itemize}
					\small
					\item Raison sociale, secteur d'activité
					\item Localisation géographique
					\item Taille / effectif
					\item Statut juridique
					\item Ancienneté
				\end{itemize}
			\end{column}
			\begin{column}{0.5\textwidth}
				\textbf{Types de besoins recueillis}
				\begin{itemize}
					\small
					\item Financement
					\item Formation / renforcement de capacités
					\item Équipement / matériel
					\item Accompagnement administratif
					\item Accès au marché
				\end{itemize}
			\end{column}
		\end{columns}
	\end{frame}
	
	% ---------------------------------------------------------------
	% PARTIE 8 : METHODOLOGIE
	% ---------------------------------------------------------------
	\section{Méthodologie d'analyse des données}
	
	\begin{frame}{Pipeline méthodologique}
		\begin{center}
			\begin{tikzpicture}[
				box/.style={draw, rounded corners, fill=institbleu!15, minimum width=2.1cm, minimum height=0.9cm, align=center, font=\scriptsize},
				>=Latex
				]
				\node[box] (boite1) at (0,0) {Nettoyage\\\& jointure};
				\node[box] (boite2) at (2.6,0) {Analyse\\univariée};
				\node[box] (boite3) at (5.2,0) {Analyse\\bivariée};
				\node[box] (boite4) at (7.8,0) {ACM};
				\node[box] (boite5) at (10.4,0) {CAH};
				\node[box] (boite6) at (13.0,0) {AFD};
				\draw[->] (boite1) -- (boite2);
				\draw[->] (boite2) -- (boite3);
				\draw[->] (boite3) -- (boite4);
				\draw[->] (boite4) -- (boite5);
				\draw[->] (boite5) -- (boite6);
			\end{tikzpicture}
		\end{center}
		\vspace{4pt}
		\small
		\begin{itemize}
			\item Nettoyage des données, sélection des variables pertinentes, jointure de \texttt{identifications.csv} et \texttt{besoins.csv}.
		\end{itemize}
	\end{frame}
	
	\begin{frame}{Analyses univariée et bivariée}
		\small
		\textbf{Analyse univariée} : décrit chaque variable isolément (fréquences, répartitions) afin de dresser un premier profil des entreprises et de leurs besoins.\\[6pt]
		\textbf{Analyse bivariée} : étudie l'association entre deux variables qualitatives à l'aide :
		\begin{itemize}
			\item du \textbf{test du $\chi^2$ d'indépendance} :
			$$\chi^2 = \sum_{i,j} \frac{(n_{ij} - e_{ij})^2}{e_{ij}}$$
			\item du \textbf{V de Cramér}, mesurant l'intensité de la liaison :
			$$V = \sqrt{\dfrac{\chi^2}{n \cdot \min(k-1, l-1)}}$$
		\end{itemize}
		\footnotesize Conditions d'interprétation : effectifs théoriques suffisants ($e_{ij} \geq 5$), $V \in [0,1]$ (proche de 0 : liaison faible ; proche de 1 : liaison forte).
	\end{frame}
	
	\begin{frame}{Analyse des Correspondances Multiples (ACM)}
		\small
		\textbf{Objectif} : synthétiser les associations entre plusieurs variables qualitatives (profils d'entreprises et besoins) dans un espace factoriel de dimension réduite.\\[6pt]
		Les valeurs propres $\lambda_k$ sont obtenues par diagonalisation de la matrice de Burt / du tableau disjonctif complet :
		$$ \mathbf{Z}^\top \mathbf{D}_n \mathbf{Z}\, u_k = \lambda_k u_k $$
		\textbf{Choix des axes} : critère des \textbf{valeurs propres corrigées de Benzécri}, qui pondère la part d'inertie réellement portée par chaque axe et corrige la sous-estimation classique de l'inertie en ACM.
	\end{frame}
	
	\begin{frame}{Classification Ascendante Hiérarchique (CAH)}
		\small
		\textbf{Objectif} : regrouper les entreprises en classes homogènes à partir de leurs coordonnées factorielles (issues de l'ACM).\\[6pt]
		\textbf{Principe} : agrégation itérative des individus/groupes les plus proches, selon un critère de distance (ex. critère de Ward minimisant l'inertie intra-classe) :
		$$ \Delta(A,B) = \frac{n_A n_B}{n_A+n_B}\, \lVert \bar{x}_A - \bar{x}_B \rVert^2 $$
		Le résultat est représenté sous forme de \textbf{dendrogramme}, dont la coupe fixe le nombre de classes retenues.
	\end{frame}
	
	\begin{frame}{Analyse Factorielle Discriminante (AFD)}
		\small
		\textbf{Objectif} : valider et modéliser l'appartenance des entreprises aux classes issues de la CAH, à partir des variables explicatives.\\[6pt]
		\textbf{Principe} : recherche des combinaisons linéaires des variables maximisant la variance inter-classes relativement à la variance intra-classe :
		$$ \lambda = \frac{\det(\mathbf{B})}{\det(\mathbf{B}+\mathbf{W})} \quad \text{(Lambda de Wilks)} $$
		Test de significativité globale par \texttt{manova()} (R), suivi d'une classification par \texttt{MASS::lda()}.
	\end{frame}
	
	% ---------------------------------------------------------------
	% PARTIE 9 : JUSTIFICATION DE LA METHODE
	% ---------------------------------------------------------------
	\section{Justification de la méthode}
	
	\begin{frame}{Pourquoi ces méthodes ?}
		\small
		Les analyses uni et bivariées sont indispensables : elles permettent un \textbf{premier diagnostic} des données, détectent les anomalies et orientent le choix des variables actives pour les analyses factorielles.
		\vspace{6pt}
		\scriptsize
		\begin{tabularx}{\textwidth}{>{\raggedright\arraybackslash}p{3.2cm} X X}
			\toprule
			\textbf{Approche} & \textbf{Principe} & \textbf{Pertinence pour ce travail} \\
			\midrule
			Régression logistique multinomiale & Modélise directement une variable cible catégorielle & Nécessite une cible pré-définie, peu adaptée à la découverte de profils \\
			Clustering k-means & Partitionnement basé sur des distances euclidiennes & Peu adapté à des données majoritairement qualitatives \\
			\textbf{ACM + CAH + AFD (retenue)} & Réduction dimensionnelle, classification puis validation discriminante & Adaptée aux données qualitatives, permet de découvrir \emph{et} valider des profils d'entreprises \\
			\bottomrule
		\end{tabularx}
	\end{frame}
	
	% ---------------------------------------------------------------
	% PARTIE 10 : RESULTATS DES ANALYSES
	% ---------------------------------------------------------------
	\section{Résultats des analyses statistiques}
	
	\begin{frame}{Résultats -- Analyse univariée}
		\small
		\textit{[Graphiques de répartition des entreprises par secteur, taille, localisation et type de besoin -- générés via la fonction \texttt{graphique\_repartition()}]}
		\begin{itemize}
			\item Prédominance de certains secteurs d'activité et types de besoins à commenter selon les résultats obtenus.
			\item Répartition géographique des entreprises enquêtées.
		\end{itemize}
	\end{frame}
	
	\begin{frame}{Résultats -- Analyse bivariée}
		\small
		\begin{itemize}
			\item Associations significatives ($\chi^2$, $p < 0.05$) entre \emph{secteur d'activité} et \emph{type de besoin dominant}.
			\item Intensité des liaisons mesurée par le V de Cramér, avec interprétation des couples de variables les plus fortement associés.
		\end{itemize}
	\end{frame}
	
	\begin{frame}{Résultats -- ACM}
		\small
		\begin{itemize}
			\item Sélection des premiers axes factoriels selon le critère de Benzécri.
			\item Interprétation des axes en fonction des modalités contributives (profil des entreprises selon leurs besoins).
		\end{itemize}
	\end{frame}
	
	\begin{frame}{Résultats -- CAH}
		\small
		\begin{itemize}
			\item Dendrogramme permettant d'identifier un nombre de classes pertinent.
			\item Caractérisation des groupes d'entreprises obtenus (profils-types de besoins).
		\end{itemize}
	\end{frame}
	
	\begin{frame}{Résultats -- AFD}
		\small
		\begin{itemize}
			\item Test de Wilks significatif, validant la séparation des groupes.
			\item \textbf{Taux de bonne classification : 97{,}24\,\%}
			\item \textbf{Coefficient Kappa : 0{,}9435} -- excellente concordance
		\end{itemize}
	\end{frame}
	
	% ---------------------------------------------------------------
	% PARTIE 11 : PLATEFORME - FONCTIONNALITES
	% ---------------------------------------------------------------
	\section{La plateforme numérique}
	
	\begin{frame}{Fonctionnalités de la plateforme}
		\begin{columns}[T]
			\begin{column}{0.5\textwidth}
				\textbf{Fonctionnalités fonctionnelles}
				\begin{itemize}
					\small
					\item Inscription et authentification des entreprises
					\item Déclaration des besoins (onboarding)
					\item Classification automatique du profil
					\item Recommandation / mise en relation ciblée
					\item Affichage des coordonnées du partenaire recommandé
				\end{itemize}
			\end{column}
			\begin{column}{0.5\textwidth}
				\textbf{Exigences non fonctionnelles}
				\begin{itemize}
					\small
					\item Sécurité des données des entreprises
					\item Performance et temps de réponse
					\item Ergonomie et accessibilité
					\item Évolutivité de la plateforme
					\item Conformité institutionnelle (charte visuelle)
				\end{itemize}
			\end{column}
		\end{columns}
	\end{frame}
	
	% ---------------------------------------------------------------
	% PARTIE 12 : OUTILS ET MODELISATION DIMENSIONNELLE
	% ---------------------------------------------------------------
	\section{Outils et modélisation dimensionnelle}
	
	\begin{frame}{Outils techniques utilisés}
		\small
		\begin{itemize}
			\item \textbf{Backend} : Django (Python), ORM Django (\texttt{annotate()}, \texttt{Count()}, \texttt{Q()})
			\item \textbf{Analyse statistique} : Python, R (\texttt{MASS::lda()}, \texttt{manova()})
			\item \textbf{Frontend} : Bootstrap, charte graphique institutionnelle (bleu marine, icônes Bootstrap Icons)
			\item \textbf{Documentation} : LaTeX (Beamer, TikZ)
		\end{itemize}
		\vspace{8pt}
		\textbf{Choix de la modélisation dimensionnelle} : approche en \textbf{schéma en étoile} inspirée des travaux de \textbf{Kimball}, considéré comme le père fondateur de la modélisation dimensionnelle, retenue pour :
		\begin{itemize}
			\small
			\item sa capacité à modéliser efficacement les faits (besoins/offres exprimés) autour de dimensions (entreprises, services) ;
			\item sa simplicité de requêtage pour l'algorithme de mise en relation.
		\end{itemize}
	\end{frame}
	
	% ---------------------------------------------------------------
	% PARTIE 13 : MODELE DE DONNEES
	% ---------------------------------------------------------------
	\section{Modèle de la base de données}
	
	\begin{frame}{Schéma en étoile et particularité de la table \texttt{Service}}
		\begin{center}
			\begin{tikzpicture}[scale=0.85,
				dim/.style={draw, rounded corners, fill=institbleu!15, minimum width=2.6cm, minimum height=0.9cm, align=center, font=\scriptsize},
				fact/.style={draw, rounded corners, fill=institor!40, minimum width=3.4cm, minimum height=1.1cm, align=center, font=\scriptsize}
				]
				\node[dim] (ent) at (-3.5,1.5) {Dimension\\\textbf{Entreprise}};
				\node[dim] (serv) at (3.5,1.5) {Dimension\\\textbf{Service}\\(offres \emph{et} besoins)};
				\node[fact] (fait) at (0,0) {Fait\\\textbf{BesoinOffreEntreprise}};
				\draw[-] (ent) -- (fait);
				\draw[-] (serv) -- (fait);
			\end{tikzpicture}
		\end{center}
		\small
		\textbf{Particularité} : la table \texttt{Service} représente simultanément les \emph{offres} et les \emph{besoins} -- un paradigme comparable au \textbf{chat de Schrödinger}, à la fois « mort et vivant » selon le contexte d'usage de l'enregistrement.\\[4pt]
		\textbf{Requête de mise en relation} : exploite l'\textbf{intersection} entre offres et besoins, facilitée par l'unicité de cette table (une seule jointure nécessaire, cf. ORM Django \texttt{annotate()}/\texttt{Count()}/\texttt{Q()}).
	\end{frame}
	
	% ---------------------------------------------------------------
	% PARTIE 14 : RESULTATS DE LA PLATEFORME
	% ---------------------------------------------------------------
	\section{Résultats de la plateforme}
	
	\begin{frame}{Résultats obtenus}
		\small
		\textit{[Captures d'écran de la plateforme : inscription, déclaration des besoins, tableau de bord des recommandations]}
		\begin{itemize}
			\item Parcours utilisateur complet : inscription $\rightarrow$ déclaration des besoins $\rightarrow$ classification $\rightarrow$ recommandations.
			\item Mise en relation effective de plusieurs profils d'entreprises testées.
			\item Interface conforme à la charte graphique institutionnelle (validée après retours sur le rendu initial).
		\end{itemize}
	\end{frame}
	
	% ---------------------------------------------------------------
	% PARTIE 15 : PERSPECTIVES
	% ---------------------------------------------------------------
	\section{Perspectives}
	
	\begin{frame}{Perspectives d'amélioration}
		\begin{itemize}
			\item Permettre aux entreprises de \textbf{se contacter directement} depuis la plateforme (messagerie intégrée).
			\item Enrichir le moteur de recommandation par des techniques de \textit{machine learning} supervisé.
			\item Élargir la couverture géographique et sectorielle de l'enquête.
			\item Mettre en place un suivi longitudinal des besoins déclarés.
			\item Intégrer des indicateurs de suivi de la SND30 directement dans la plateforme.
		\end{itemize}
	\end{frame}
	
	% ---------------------------------------------------------------
	% CONCLUSION
	% ---------------------------------------------------------------
	\section*{Conclusion}
	\begin{frame}{Conclusion}
		\small
		Ce travail a permis de :
		\begin{itemize}
			\item cartographier statistiquement les besoins des entreprises camerounaises à partir de l'enquête MINEPAT 2025 ;
			\item concevoir et développer une plateforme opérationnelle de mise en relation, fondée sur un modèle dimensionnel simple et efficace ;
			\item valider la robustesse de la démarche méthodologique (ACM, CAH, AFD) par un taux de classification de 97{,}24\,\% et un Kappa de 0{,}9435.
		\end{itemize}
		\vspace{10pt}
		\centering \large \textbf{Merci de votre attention}
	\end{frame}
	
	\begin{frame}
		\centering
		\Huge \textbf{Questions ?}
	\end{frame}
	
\end{document}
